% Right-divergence and constrained-current members; no preamble.
% Requires amsmath/amssymb/amsthm, theorem environment, and \norm.
% Depends on the preceding pt- and pc- results. Labels have prefix rd-.
\section{Right-divergence data and the constrained current inverse}
\label{rd-section}

The remaining right-divergence estimates require cancellation on
H\"older input. We isolate that cancellation in a translation-invariant
massive lattice Green function. The argument applies to the full
homogeneous flat $G_0$ and its auxiliary and full current perturbations.
It does not identify the distinct scalar inverse $G'$ with $G_0$.

Use the admissible spatial weights and weak norms of the preceding
section. Let $\norm{\lambda}_{\gamma,\varpi}$ denote the weighted
supremum of the fine field and its $\gamma$-H\"older seminorm on
fixed cell collars. Fields are kept complete; no H\"older datum is
multiplied by a sharp cell indicator.

\begin{theorem}\label{rd-theorem}
For $0\le\beta<\gamma<1$ and a sufficiently small fixed spatial
exponent $\kappa$, the homogeneous flat inverses satisfy
\begin{equation}\label{rd-main-estimate}
 N_{\beta,\varpi}(G D^*\lambda)
       \le C_{\beta,\gamma}\norm{\lambda}_{\gamma,\varpi},
 \qquad G\in\{G_0,G_\pi^J,G_1^J\}.
\end{equation}
The current lies in the previously established common pointwise and
Hilbert ball. Constants are independent of refinement and compatible
periodic volume. Here $D^*$ can be any specified physical fine
divergence, with its fixed orientation contractions. Taking
$\gamma=\beta+\varepsilon$, with fixed
$0<\varepsilon<1-\beta_0$, makes the constants uniform for
$0\le\beta\le\beta_0$. The source choice
$\varepsilon=(1-\beta_0)/2$ is included.
\end{theorem}

This is a bound on the right-differentiated operator, not an estimate
obtained by bounding $\norm{D^*\lambda}_\infty$. The constant can
grow as $\gamma\downarrow\beta$; the endpoint is not claimed.
The relevant source members are described at deposited Raw I,
lines 13601--13615 and 13759--13765.

\subsection{A remainder that accepts arbitrary bounded data}

Introduce the comparison operator
\begin{equation}\label{rd-massive-split}
 G_m=(\mathcal L+1)^{-1}I,\qquad R_m=G_0-G_m.
\end{equation}
It is an auxiliary decomposition, not a change of source precision.
Shrink the common complex chart so that $P_0+1$ is uniformly nonzero;
on the real chart it is at least one. The high denominators already
have this property. The high scalar diagonal of $R_m$ is
\begin{equation}\label{rd-high-remainder}
 \frac1{P_r}-\frac1{P_r+1}
    =\int_0^\infty(1-e^{-t})e^{-tP_r}\,dt,
                    \qquad r\ne0.
\end{equation}
The other terms are the finite-rank corrections of $G_0$ and the
bounded central subtraction $-(P_0+1)^{-1}\Pi_0$.

With one input derivative, one output derivative, and an output
$\beta$-H\"older increment, the small-time bound from the heat
estimates is
\[
 (1-e^{-t})t^{-(2+\beta)/2}\le t^{-\beta/2},
                       \qquad0<t\le1.
\]
It is integrable. Lower output members are easier, and large times
use the high-alias gap. In each rank term, the input derivative
multiplies its row by one fine multiplier. The high row
$\omega_s/P_s$ becomes $O(|\omega_s|/(1+|s|))$; the Woodbury row
is bounded by
$C(|a_s|/(1+|s|)^2+|\omega_s|/(1+|s|))$.
Their alias $\ell^2$ norms are uniform. The output
$C^{1,\beta}$ vectors retain the proved absolute alias sums for
$\beta<1$.

The full reconstructed remainder is periodic and holomorphic in
the resulting pointwise operator norms. Chart gluing and contour
shifting therefore give exponential coarse-cell kernels and
\begin{equation}\label{rd-remainder-bound}
 N_{\beta,\varpi}(R_mD^*f)
                  \le C_\beta\norm{f}_{\infty,\varpi},
                       \qquad0\le\beta<1.
\end{equation}
This remainder estimate genuinely accepts arbitrary bounded fine
data. It requires no H\"older extension across a cell boundary:
the right derivative has already been placed on its kernel.
Compatible periodic sampling gives the same bound on tori.

\subsection{Weighted heat moments down to the mesh scale}

Let $p_t^\eta(v)$ be the full-lattice probability heat kernel of
$\mathcal L$, with fine displacement $v\in\eta\mathbb Z^d$.
For $0<\gamma<1$ and $k=2,3$, we have
\begin{equation}\label{rd-heat-moment}
 \sum_v e^{\kappa|v|_1}|v|_1^\gamma
                   |D_\eta^kp_t^\eta(v)|
       \le Ct^{(\gamma-k)/2},\qquad0<t\le1.
\end{equation}
The probability-kernel convention includes the physical normalization:
if a kernel density is used instead, its $\eta^{-d}$ factor is
canceled by physical summation.

For $t\ge\eta^2$, apply the weighted heat variation estimate with
physical exponential scale $\delta=\kappa+t^{-1/2}$. Its lattice
parameter $\eta\delta$ lies in a fixed bounded interval. Semigroup
splitting into $k$ factors gives
\[
 \sum_v e^{\delta|v|_1}|D_\eta^kp_t^\eta(v)|
                \le C_kt^{-k/2},\qquad \eta^2\le t\le1,
\]
because $\delta^2t$ is bounded. The inequality
$|v|^\gamma\le C_\gamma t^{\gamma/2}e^{|v|/\sqrt t}$
proves \eqref{rd-heat-moment} in this range.

For $0<t<\eta^2$, expand the finite difference into its fixed list
of translates:
\[
 D_\eta^kp_t^\eta(v)
    =\eta^{-k}\sum_jc_jp_t^\eta(v+\eta a_j),
 \qquad |a_j|_1\le k,\quad\sum_j|c_j|\le C_k.
\]
The dimensionless random walk has time $t/\eta^2<1$, so its
exponential and fixed positive moments are uniformly bounded.
The weighted moment above is at most $C\eta^{\gamma-k}$.
Since $k>\gamma$ and $t<\eta^2$,
$\eta^{\gamma-k}\le t^{(\gamma-k)/2}$, proving
\eqref{rd-heat-moment} also below the mesh time.

For $t\ge1$, absorb $|v|^\gamma$ into a fixed additional
exponential weight $e^{\delta_1|v|}$. The weighted heat estimates
are bounded by $Ce^{c(\kappa+\delta_1)^2t}$. Choose
$\kappa,\delta_1$ small enough that this is at most $Ce^{t/2}$.
The massive factor $e^{-t}$ leaves an integrable tail. The required
zeroth and first heat derivatives obey the analogous estimates
without the polynomial moment.

\subsection{Cancellation on complete fields}

Every nonzero difference of the translation-invariant heat kernel
has exact zero row sum:
\begin{equation}\label{rd-zero-row}
 \sum_vD_\eta^kp_t^\eta(v)=0\qquad(k\ge1).
\end{equation}
The local weighted H\"older norm implies
\begin{equation}\label{rd-global-increment}
 \varpi(\operatorname{cell}x)|\lambda(y)-\lambda(x)|
 \le C\norm{\lambda}_{\gamma,\varpi}
        e^{\kappa|x-y|_1}|x-y|_1^\gamma.
\end{equation}
At bounded separation use a common fixed collar and neighboring
weight ratios. At larger separation use the two supremum bounds;
the polynomial factor is then bounded below and the cell-weight
ratio is controlled by physical distance up to a fixed offset.
No derivative of the weight is taken.

Subtract $\lambda(x)$ using \eqref{rd-zero-row}. Equations
\eqref{rd-heat-moment}--\eqref{rd-global-increment} give
\begin{align}
 \norm{D_\eta^2e^{-t\mathcal L}\lambda}_{\infty,\varpi}
   &\le Ct^{-1+\gamma/2}\norm{\lambda}_{\gamma,\varpi},
       \label{rd-heat-second}\\
 \norm{D_\eta^3e^{-t\mathcal L}\lambda}_{\infty,\varpi}
   &\le Ct^{-3/2+\gamma/2}\norm{\lambda}_{\gamma,\varpi},
                   \qquad0<t\le1.
       \label{rd-heat-third}
\end{align}
The difference between two second-derivative output values is
bounded by twice \eqref{rd-heat-second}, and by their physical
separation times \eqref{rd-heat-third}. Interpolating yields
\begin{equation}\label{rd-heat-holder}
 [D_\eta^2e^{-t\mathcal L}\lambda]_{\beta,\varpi}
       \le Ct^{-1+(\gamma-\beta)/2}
                                  \norm{\lambda}_{\gamma,\varpi}.
\end{equation}
Its integral is finite, with
\begin{equation}\label{rd-time-integral}
 \int_0^1t^{-1+(\gamma-\beta)/2}\,dt
                      =\frac2{\gamma-\beta}.
\end{equation}
For $\beta=0$ the supremum estimate suffices. The field member of
$G_mD^*\lambda$ uses only the first heat derivative, whose
$t^{-1/2}$ bound is integrable on bounded input. The gradient and
its H\"older member use
\eqref{rd-heat-second}--\eqref{rd-time-integral}. The large-time
massive tail is integrable as above. Therefore
\begin{equation}\label{rd-massive-bound}
 N_{\beta,\varpi}(G_mD^*\lambda)
                  \le C_{\beta,\gamma}
                                  \norm{\lambda}_{\gamma,\varpi}.
\end{equation}
All singular-integral cancellation has occurred in this global
translation-invariant calculation. On tori use the periodized heat
kernel, or periodically extend the fields and weights to the cover.
Periodic distance is no larger than cover distance, so the same
moment bounds and multiplicities give identical constants.

\subsection{Transfer to the actual current inverses}

Adding \eqref{rd-remainder-bound} and \eqref{rd-massive-bound}
proves \eqref{rd-main-estimate} for $G_0$. For either current
family, the proved weak-space perturbation obeys
\[
 N_{\beta,\varpi}(K_JA)
       \le C_*\norm{J}_\infty N_\varpi(A),\qquad
 C_*\norm{J}_\infty<\tfrac12.
\]
Apply the same convergent series to the baseline datum:
\begin{equation}\label{rd-current-transfer}
 G^JD^*\lambda=(I-K_J)^{-1}G_0D^*\lambda.
\end{equation}
The baseline $N$ bound requires positive input H\"older regularity;
the baseline $N_\beta$ bound uses $\gamma>\beta$. Every later term
is controlled by the established $N$-to-$N_\beta$ perturbation
bound. This proves the remaining assertions of
\eqref{rd-main-estimate}. The larger H\"older-data constant does
not shrink the current radius: it multiplies the initial datum,
not the contraction. Holomorphy follows from the same Neumann series.

Localization is legitimate with a smooth physical cutoff $\phi$
of fixed support collar and bounded first derivative. The product
difference identity gives
$\norm{\phi\lambda}_{\gamma,\varpi}
 \le C_\phi\norm{\lambda}_{\gamma,\varpi;\,	ext{collar}}$
without a power of $\eta^{-1}$. A source-centered exponential weight
then gives decay away from the support. A sharp cell indicator
would generally create $\eta^{-\gamma}$ H\"older jumps and is not
permitted by this argument. This matches the use of localized
cutoffs in the stated source member.

\subsection{The constrained-current member}

The source separately defines, at Raw I, line 13735,
\begin{equation}\label{rd-constrained-definition}
 \widetilde G^J=(I-\Gamma\Theta_J)^{-1}\Gamma,
 \qquad
 \Gamma=G_0-H_{\rm src}q_{\rm aux}^{-1}H_{\rm src}^*.
\end{equation}
The preceding section proved $Q\Gamma=0$ and, for a uniform
constant $B_\Gamma$,
$N_{\beta,\varpi}(\Gamma f)
 \le B_\Gamma\norm{f}_{\infty,\varpi}$.
It also proved
$\norm{\Theta_Jf}_{\infty,\varpi}
 \le\tau_\infty C_v\norm{J}_\infty
                      \norm{f}_{\infty,\varpi}$.
Decrease the common current radius to
\begin{equation}\label{rd-constrained-radius}
 \varepsilon_{\rm all}
  =\min\left\{\varepsilon_{\rm pt},
             \frac1{2B_\Gamma\tau_\infty C_v}\right\},
\end{equation}
omitting the second restriction if its coefficient is zero.
The series in \eqref{rd-constrained-definition} then converges in
weighted pointwise norms and gives the source (J6) estimate
\begin{equation}\label{rd-constrained-bound}
 N_{\beta,\varpi}(\widetilde G^Jf)
           \le2B_\Gamma\norm{f}_{\infty,\varpi},\qquad
                 Q\widetilde G^J=0.
\end{equation}
For the bound, estimate the outside $\Gamma$ in $N_\beta$ and
sum the remaining supremum-norm geometric series. For the exact
constraint, every term begins on the left with $\Gamma$, so applying
$Q$ gives zero term by term; the limit preserves the identity.
The map is analytic, and source-centered weights give exponential
coarse-cell locality. This statement keeps the source's constrained
inverse distinct from $G_\pi^J$ and $G_1^J$.

\subsection{Remaining distinctions}

The right-divergence theorem covers $G_0,G_\pi^J,G_1^J$ for the
stated H\"older input and also supplies the separate constrained
current member \eqref{rd-constrained-bound}. The deposited source
also names a right-divergence estimate for the distinct scalar
inverse $G'$. That member is not inferred by replacing $G'$ with
$G_0$; it requires its own scalar Woodbury adaptation and the
corresponding scalar penalty normalization. It remains separate here.

No current Laplacian estimate or equal-exponent endpoint
$\gamma=\beta$ is asserted. The results are homogeneous and flat,
with the corrected fixed-step constraint term already matched to
the deposited transcription. They do not prove nonflat source
stability, mixed retained-domain bounds, strict background locality,
or an interacting Yang--Mills ultraviolet limit.
